\section{Detector Error Model}

In this section we want to determine the error model which we use to construct the decoders. 

For that we determine the probability with which the detectors detects certain errors. 
We call this the \enquote{detection error model}.

\begin{enumerate}
    \item \textcolor{red}{Put here a nice description, how we determine the difference between the data qubit and the ancilla qubit as measurement error.}
    \item \textcolor{red}{Show circuit with error locations}
    \item \textcolor{red}{Describe the whole appraoch}
\end{enumerate}



\subsection{General Statements}
We first discuss all the general statements that hold for all Pauli channels before showing how each individual pauli channel behaves. 


\subsubsection{Error Channels of our Noise Model}

Our circuit level noise model consists of the following Pauli error channels:

Bit-flip error channel:
\begin{equation}
    N_{X}(\rho) = (1-p) \rho + p \cdot X \rho X
\end{equation}
which we use to simulate measurement errors. 
For this we place those error channels before the ancialla qubit measurement.


Depolarizing one qubit channel:
\begin{equation}
    N_{D}(\rho) = (1-p) \rho + \frac{p}{3} \sum^3_{i=1} E_i \rho E_i
\end{equation}
\begin{equation*}
    E_i \in \{ X, Y, Z \}
\end{equation*}
which are placed after the initialization of the logical qubits, to mimic a fault tolerant prepartion procedure.

Depolarizing two qubit channel
\begin{equation}
    N_{D_2}(\rho) = (1-p) \rho + \frac{p}{15} \sum_{i=1}^{15} E_i \rho E_i,
\end{equation}
\begin{equation*}
    E_i \in \{E_k \otimes E_l | E_k,E_l \in \{I,X,Y,Z\} \} \setminus \{I\otimes  I\}
\end{equation*}
which are placed after each CNOT to mimic simulate a gate errors.


\subsubsection{Pauli Channel Commute}

It is easy to show that all Pauli channels commute.
We first  define two arb. Pauli channels.\footnote{Complex conjugate of a Pauli is the Pauli ${P}^\dag={P}$}

\begin{equation}
    N_1 (\rho) = \sum_P p_P {P} \rho {P}
\end{equation}

\begin{equation}
    N_2 (\rho) = \sum_Q p_Q {Q} \rho {Q}
\end{equation}

We also know that all Paulis either commute or anticommute and therefore we can write
\begin{equation}
    {P}{Q} = (-1)^\sigma {Q} {P},
\end{equation}
where $\sigma$ is either $0$ if they commute or $1$ if they anticommute.

If we look at the composition of the two arb. Pauli channels, using the above identity for pauli channels
\begin{align*}
    N_2  (N_1 (\rho)) & = \sum_{P,Q} p_P p_Q {Q} {P} \rho {P} {Q} \\
    & = \sum_{P,Q} p_P\  p_Q (-1)^\sigma {P} {Q} \rho (-1)^\sigma {Q} {P} \\
    & = (-1)^{2\cdot\sigma} \sum_{P,Q} p_P p_Q  {P} {Q} \rho  {Q} {P} \\
    & = (1)^\sigma \sum_{P,Q} p_P p_Q  {P} {Q} \rho  {Q} {P}
    = N_1  (N_2 (\rho))
\end{align*}
we can see that they commute.


\subsubsection{Pauli Operator and CNOTs}

To move the error channels through the circuit we need to move them through the CNOT gates.
Therefore we state here the basic CNOTs error propagation identities. 

$X$-Errors pass trivial through the NOT part of the CNOT and a propagted to the other qubit by the control side.
And $Z$-Errors behave the opposite way.

\subsubsection{Stabilizer Equivalence}

We can define stabilizer equivalences ($\tilde{s}$) based on the measurement we perform to determine our stabilizer.
For example if we measure our stabilizer in the $Z$-basis, 
phase-flips will not change the stabilizer measurement and therefore we can group them together with the identity.
By the same idea both $Y$- and $X$-errors lead to the same syndrome and can be group as such together in the sense of stabilizer equivalence. 

For $Z$-measurements:
\begin{align*}
    X & \overset{\tilde{s}}{\rightarrow} X \\
    Y & \overset{\tilde{s}}{\rightarrow} X \\
    Z & \overset{\tilde{s}}{\rightarrow} I \\
    I & \overset{\tilde{s}}{\rightarrow} I \\
\end{align*}

For $X$-measurements:
\begin{align*}
    Z & \overset{\tilde{s}}{\rightarrow} Z \\
    Y & \overset{\tilde{s}}{\rightarrow} Z \\
    X & \overset{\tilde{s}}{\rightarrow} I \\
    I & \overset{\tilde{s}}{\rightarrow} I \\
\end{align*}

A depolarizing channel which is infront of a $Z$-measurement can be therefore rewritten as
\begin{align*}
    N_D(\rho) = & \left( 1- \frac{4}{3} p_D\right) \hat{I} + \frac{1}{3} p_D \sum_{i=0}^3 \hat{E_i} \\
\overset{\tilde{s}}{\rightarrow} & \left(  1- \frac{2}{3}p_D\right) + \frac{2}{3} p_D \hat{X}
\end{align*}
a bit flip channel with probability $p_x = 2/3 p_D$.


\subsubsection{Compostion of Channels (Serial Combination)}

In this section we find general formulars for the composition of single qubit depolarizing channels (depo) and bit- and phase-flip channels.
These are the only channels which will appear after commuting the channels to the end of the circuit and tracing out the data qubit
and therefore the only one of interest. 

To make the calculations shorter we rewrite the depo channel, 
by including the identity as $E_0=I$ in the sum, s.t. 
\begin{align*}
    N_{D1}(\rho)    &= (1-p) \rho  + \frac{1}{3} p \sum_{i=1}^3 E_i \rho E_i \\
                    &= (1-\frac{4}{3}p)\rho + \frac{1}{3} p \sum_{i=0}^{3} E_i \rho E_i.
\end{align*}
Note that the summation starts at $i=0$ and therefore 
\begin{equation*}
    E_i \in \{ I, X, Y, Z\}. 
\end{equation*}


Furthermore we to save on space we rewrite the channels as a mixture\footnote{
    \textcolor{red}{Unsure if this is the correct term.}} 
of (Pauli-)channels using the notation

\begin{equation*}
    \hat{E_i} \hat{=} E_i \rho E_i^\dag.
\end{equation*}


\paragraph{Depo-Depo Channel}

The serial combination of depo channels can be rewritten as 
\begin{align*}
    N_{D;p_2}(N_{D;p_1}(\rho)) &= \left( \left(1-\frac{4}{3}p_1\right) \hat{I} + \frac{1}{3} p_1 \sum_{i=0}^{3}\hat{E_i} \right)
    \cdot \left( \left(1-\frac{4}{3}p_2\right) \hat{I} + \frac{1}{3} p_2 \sum_{i=0}^{3}\hat{E_i} \right) \\
    &= \left(1-\frac{4}{3}p_1\right)\left(1-\frac{4}{3}p_2\right) \hat{I} \\
    &  \ \ \ \ + \frac{1}{3}\left((1-\frac{4}{3}p_1)p_2 + (1-\frac{4}{3}p_2)p_1 + \frac{4}{3}p_1 p_2\right) \sum_{i=0}^{3} \hat{E_i} \\
    &= \left( 1- \frac{4}{3}\left( p_1 + p_2 - \frac{4}{3}p_1 p_2\right) \right) \hat{I} 
    + \frac{1}{3} \left( p_1 + p_2 - \frac{4}{3}p_1 p_2\right) \sum_{i=0}^{3} \hat{E_i} \\
    &= (1-\frac{4}{3}p_c)\hat{I} + \frac{1}{3} p_c \sum_{i=0}^{3} \hat{E_i}\\
    &= N_{D;p_c}(\rho)
\end{align*}

Because we include the identity operator in the sum, the product of the sums becomes a trivial cyclic/symmetric remapping 
\begin{equation*}
    \left( \sum_{i=0}^{3} \hat{E_i}\right) \cdot \left( \sum_{i=0}^{3} \hat{E_i}\right) = 4 \left( \sum_{i=0}^{3} \hat{E_i}\right) .
\end{equation*}

From above calculation we can conclude that the composition of two depolarizing one qubit channels 
is just a depolarizing one qubit channel with the combined probability 
\begin{equation*}
    p_c = f_{2}(p_1,p_2) = p_1 + p_2 - \frac{4}{3}p_1 p_2 = p_1 + (1- \frac{4}{3}p_1) p_2
\end{equation*}

Note that the combined probability is indepent of the ordering of the channels (as expected).

Using the function for two probabilities $f_{2}$ we can define a recursive function $f_{n}$ to determine the combine probability of the composition of multiple depolarizing channels.
\begin{align*}
    f_{n}(p_1,p_2,p_3,...,p_n) &= f_{2}\left(p_1,f_{n-1}(p_2,p_3,...,p_n)\right) \\
    &= f_{2}(p_1,f_{2}(p_2,f_{n-2}(p_3,...,p_n)))
\end{align*} 


\paragraph{Flip-Flip Channel}

An even number of flips cancels each other and an odd number of flips leads to a resulting flip error.
Therefore the composition of multiple (bit-/phase-)flip channels is once again a (bit-/phase-)flip channel with the combined probability
\begin{equation*}
    p_c = \sum_{\vec{n}\in\{\vec{n} \text{ odd}\}} \prod_i p_i^{n_i} \cdot \left( 1-p_i\right)^{1-n_i}
\end{equation*}
where we sum over all possible odd occupation vectors $\vec{n}$ of errors.\footnote{The occupation vector is defined s.t. if $n_i = 1$ the $i$-th flip channel flips the bit}.
An occupation vector is defined as odd if 
\begin{equation*}
    \left(\sum_i n_i\right) \mod 2 = 1.
\end{equation*}




\paragraph{Depo-(Bit-)Flip Channel}

The composition of a (bit-)flip and depo channel takes the form
\begin{align*}
    N_{X}(N_{D}(\rho)) &= \left( (1-p_x) \hat{I} + p_x \hat{X}\right)
    \cdot \left( \left(1-\frac{4}{3}p_D\right) \hat{I} + \frac{1}{3} p_D \sum_{i=0}^{3}\hat{E_i} \right) \\
    &= (1-p_x)\left(1-\frac{4}{3}p_D\right) \hat{I} + \frac{1}{3} p_D \sum_{i=0}^{3}\hat{E_i}  + p_x\left(1-\frac{4}{3}p_D\right) \hat{X}
\end{align*}
where we used
\begin{equation*}
    \sum_{i=0}^{3} \hat{X}\hat{E_i} = 
    \sum_{i=0}^{3} \hat{E_i}. 
\end{equation*}


\paragraph{All Channels}

To find the composition of bit- and phase-flip and depo channel 
we combine the previous results with an phase-flip channel to get 
\begin{align*}
    N_{final} =&N_Z(N_X(N_D(\rho))) \\
    =& \left((1-p_z)\hat{I} + p_z \hat{Z}\right) \\
    & \cdot \left((1-p_x)\left(1-\frac{4}{3}p_D\right) \hat{I} + \frac{1}{3} p_D \sum_{i=0}^{3}\hat{E_i}  + p_x\left(1-\frac{4}{3}p_D\right) \hat{X}\right) \\
    =& \frac{1}{3} p_D \sum_{i=0}^{3} \hat{E_i} \\
    &+ \left(1-\frac{4}{3}p_D \right)\cdot\left((1-p_x)(1-p_z) \hat{I}+ (1-p_x) p_z \hat{Z} + p_x (1-p_z) \hat{X} + p_x p_z \hat{Y}\right)
\end{align*}

From this combination we can extract the final probabilities $p_{i}^t$ of detecting each Pauli $E_i$ 
dependend on the defining probabilities of the depo $p_D$, bit-flip $p_x$ and phase-flip $p_Z$ channels.
\begin{align}
    p_0^t = p_I^t =& \frac{1}{3} p_D + \left( 1-\frac{4}{3}p_D\right) (1-p_x)\cdot(1-p_z) \\
    p_1^t = p_x^t =& \frac{1}{3} p_D + \left( 1-\frac{4}{3}p_D\right) p_x \cdot(1-p_z) \\
    p_2^t = p_y^t =& \frac{1}{3} p_D + \left( 1-\frac{4}{3}p_D\right) p_x\cdot p_z\\
    p_3^t = p_z^t =& \frac{1}{3} p_D + \left( 1-\frac{4}{3}p_D\right) (1-p_x)\cdot p_z \\
\end{align}





